\chapter{Droites et systèmes linéaires}
\section*{Exercice 10.1 -- Deux points}
Sa pente vaut $(9-3)/(4-1)=2$, puis $p=3-2=1$. Une équation est $y=2x+1$, ou $2x-y+1=0$.

\section*{Exercice 10.2 -- Point et directeur}
Le déterminant $3(y+1)-2(x-2)=0$ donne $-2x+3y+7=0$. Le point vérifie l'équation et un directeur est bien $(3,2)$.

\section*{Exercice 10.3 -- Verticale}
Elle a pour équation $x=4$, et pour directeur $(0,1)$. La différence des abscisses est nulle : le quotient de pente est indéfini. Ce n'est pas la courbe d'une fonction de $x$.

\section*{Exercice 10.4 -- Intersection}
L'égalité $3x-2=-2x+8$ donne $5x=10$, donc $x=2$, puis $y=4$. L'intersection est $(2,4)$.

\section*{Exercice 10.5 -- Trois systèmes}
Le premier système décrit la même droite deux fois : infinité de couples $(x,3-x)$. Le deuxième est contradictoire, car doubler la première ligne imposerait $6=7$. Dans le troisième, l'addition donne $2x=4$, donc l'unique couple est $(2,1)$.

\section*{Exercice 10.6 -- Région admissible}
C'est l'intersection du premier quadrant et du demi-plan sous $y=6-2x$, frontière comprise. Le triangle plein a pour sommets $(0,0),(3,0),(0,6)$. Les intersections avec les axes se trouvent en annulant successivement $y$ et $x$.
